\documentclass[12pt, twoside]{article}
\input{/Users/chris/github/course-files/docs/ib/preamble}

\fancyhead[LE]{\thepage}
\fancyhead[RO]{Markscheme}
\fancyhead[LO]{La Scuola d'Italia / Dr.\ Huson / IB Math 12 HL \\* 18 November 2026}

\begin{document}
\subsubsection*{5.6 Classwork: The Logistic Equation (no calculator) --- Markscheme}

\emph{Total: 38 marks.}

\begin{enumerate}[itemsep=0.3cm]

%--- Problem 1: Forming differential equations from words [7 marks] ---
\item[1.] [Maximum mark: 7]

\textbf{(a)} $\dfrac{\mathrm{d}P}{\mathrm{d}t} = kP$ \hfill \textbf{A1}

\textbf{(b)} $\dfrac{\mathrm{d}A}{\mathrm{d}t} = -kA$ \hfill \textbf{A1}

\textbf{(c)} $\dfrac{\mathrm{d}V}{\mathrm{d}t} = -k\sqrt{V}$ \hfill \textbf{A1}

\emph{(b) and (c): the negative sign is required, because $k > 0$ and the
quantity decreases. $\frac{\mathrm{d}A}{\mathrm{d}t} = kA$ earns A0 (the
5-4CW 1(d) error).}

\textbf{(d)} $\dfrac{\mathrm{d}N}{\mathrm{d}t} = kN\left(400 - N\right)$
\hfill \textbf{A1A1}

\emph{A1 for a rate proportional to a product with $N$. A1 for
$400 - N$, the number who have not heard. $kN(N - 400)$ earns A1A0.}

\textbf{(e)} When $N$ is small, the factor $N$ is small: few students
know the rumour to pass it on. \hfill \textbf{R1}

When $N$ is close to $400$, the factor $400 - N$ is small: few students are
left to hear it. \hfill \textbf{R1}

\emph{Each R1 needs a reference to the factor in the equation (or the
product being small), not just the context.}

\bigskip \hrule \bigskip

%--- Problem 2: Partial fractions; solve the logistic DE; rearrange [10 marks] ---
\item[2.] [Maximum mark: 10]

\textbf{(a)} $1 = A(400 - N) + BN$ \hfill \textbf{M1}

$N = 0$: $A = \dfrac{1}{400}$; \quad $N = 400$: $B = \dfrac{1}{400}$ \hfill \textbf{A1}

\emph{A1 for both values.}

\textbf{(b)} $\displaystyle\int \frac{1}{N(400 - N)}\,\mathrm{d}N = \int k\,\mathrm{d}t$
\hfill \textbf{M1}

$\dfrac{1}{400}\left(\ln N - \ln(400 - N)\right) = kt + C$ \hfill \textbf{A1A1}

$\ln\left(\dfrac{N}{400 - N}\right) = 400kt + 400C$

$t = 0$, $N = 40$: \ $400C = \ln\dfrac{40}{360} = \ln\dfrac{1}{9} = -\ln 9$ \hfill \textbf{M1A1}

$\ln\left(\dfrac{N}{400 - N}\right) = 400kt - \ln 9$ \hfill \textbf{AG}

\emph{M1 for separating with integral signs. A1 for $\ln N$; A1 for
$-\ln(400 - N)$, with the negative from the chain rule. Moduli are not
needed, since $0 < N < 400$. M1 for using $t = 0$, $N = 40$; A1 for the
constant. Wrong method: $+\ln(400 - N)$ gives
$\ln\left(N(400 - N)\right)$ and the given answer cannot be reached: A0 for
that term and A0 for the constant.}

\textbf{(c)} $\dfrac{N}{400 - N} = e^{400kt - \ln 9} = \dfrac{e^{400kt}}{9}$
\hfill \textbf{M1}

$9N = 400e^{400kt} - Ne^{400kt}$, \quad $N\left(9 + e^{400kt}\right) = 400e^{400kt}$
\hfill \textbf{M1}

$N = \dfrac{400e^{400kt}}{9 + e^{400kt}} = \dfrac{400}{1 + 9e^{-400kt}}$
\hfill \textbf{A1}

\hfill \textbf{AG}

\emph{First M1 for exponentiating both sides correctly. Wrong method:
exponentiating term by term, $\frac{N}{400 - N} = e^{400kt} - 9$, earns M0
(the 5-4CW 2(b) error). Second M1 for collecting the $N$ terms. A1 for
dividing numerator and denominator by $e^{400kt}$, seen.}

\newpage

%--- Problem 3: Logistic solution with data [11 marks] ---
\item[3.] [Maximum mark: 11]

\textbf{(a)} $100 = \dfrac{400}{1 + 9e^{-800k}}$, \quad $1 + 9e^{-800k} = 4$
\hfill \textbf{M1}

$e^{-800k} = \dfrac{1}{3}$, \quad $800k = \ln 3$, \quad $k = \dfrac{\ln 3}{800}$
\hfill \textbf{A1}

\hfill \textbf{AG}

\textbf{(b)} $200 = \dfrac{400}{1 + 9e^{-400kt}}$, \quad $9e^{-400kt} = 1$
\hfill \textbf{M1}

$400kt = \ln 9$, \quad $t = \dfrac{\ln 9}{400k} = \dfrac{2\ln 3}{\frac{1}{2}\ln 3}$
\hfill \textbf{A1}

$t = 4$ days \hfill \textbf{A1}

\emph{$400k = \frac{\ln 3}{2}$. Wrong method: a straight line through
$(0, 40)$ and $(2, 100)$ gives $t = 5\frac{1}{3}$: M0A0A0.}

\textbf{(c)} $400k = \dfrac{\ln 3}{2}$, so
$e^{-400kt} = e^{-\frac{t}{2}\ln 3} = 3^{-t/2}$ \hfill \textbf{M1}

$9e^{-400kt} = 3^{2} \times 3^{-t/2} = 3^{2 - t/2}$ \hfill \textbf{A1}

\hfill \textbf{AG}

\textbf{(d)} $N = \dfrac{400}{1 + 3^{2 - 4}} = \dfrac{400}{1 + \frac{1}{9}}$
\hfill \textbf{M1}

$N = 360$ \hfill \textbf{A1}

\emph{Wrong method: continuing the pattern $40$, $100$, $200$ in a straight
line from $t = 2$ to $t = 4$ gives $400$ at $t = 8$: M0A0.}

\textbf{(e)} $N \to 400$ \hfill \textbf{A1}

In the long run every student in the school hears the rumour. \hfill \textbf{R1}

\emph{The R1 needs the context (all $400$ students), not just ``the
population levels off''.}

\bigskip \hrule \bigskip

%--- Problem 4: General logistic: second derivative, maximum rate [10 marks] ---
\item[4.] [Maximum mark: 10]

\textbf{(a)} $\dfrac{\mathrm{d}P}{\mathrm{d}t} = k\left(NP - P^{2}\right)$, so
$\dfrac{\mathrm{d}^{2}P}{\mathrm{d}t^{2}} = k\left(N - 2P\right)\dfrac{\mathrm{d}P}{\mathrm{d}t}$
\hfill \textbf{M1A1}

$= k\left(N - 2P\right) \cdot kP\left(N - P\right)
= k^{2}P\left(N - P\right)\left(N - 2P\right)$ \hfill \textbf{A1}

\hfill \textbf{AG}

\emph{M1 for differentiating with respect to $t$ using the chain rule
(implicit differentiation). A1 for $k(N - 2P)\frac{\mathrm{d}P}{\mathrm{d}t}$.
Differentiating with respect to $P$ only, $k(N - 2P)$, without the factor
$\frac{\mathrm{d}P}{\mathrm{d}t}$: M0.}

\textbf{(b)} $\dfrac{\mathrm{d}^{2}P}{\mathrm{d}t^{2}} = 0$, and $P \neq 0$,
$P \neq N$, so $P = \dfrac{N}{2}$ \hfill \textbf{M1A1}

For $P < \frac{N}{2}$, $\frac{\mathrm{d}^{2}P}{\mathrm{d}t^{2}} > 0$; for
$P > \frac{N}{2}$, $\frac{\mathrm{d}^{2}P}{\mathrm{d}t^{2}} < 0$. So
$\frac{\mathrm{d}P}{\mathrm{d}t}$ changes from increasing to decreasing, and
is greatest at $P = \frac{N}{2}$. \hfill \textbf{R1}

\emph{Accept: $\frac{\mathrm{d}P}{\mathrm{d}t} = k(NP - P^{2})$ is a
quadratic in $P$ with a negative $P^{2}$ coefficient, so its maximum is at
the vertex $P = \frac{N}{2}$ (M1A1R1). The R1 needs a reason, not just
``second derivative is zero'' (the 5-4CW 4(c) error).}

\textbf{(c)} $k \cdot \dfrac{N}{2}\left(N - \dfrac{N}{2}\right)$ \hfill \textbf{M1}

$= \dfrac{kN^{2}}{4}$ \hfill \textbf{A1}

\textbf{(d)} $\dfrac{\ln 3}{800} \times \dfrac{400^{2}}{4} = 50\ln 3$ students
per day \hfill \textbf{A1}

at $t = 4$ days, when $N = 200$ \hfill \textbf{A1}

\emph{FT from (c) and from 3(a), 3(b). Accept $\frac{200\ln 3}{4}$ or any
exact equivalent.}

\end{enumerate}
\end{document}
